\section{模的一般性质}\label{sec:galois-1-3}
模对于（交换）环，正如向量空间对于域. 不过，模未必有基，这使得模的研究困难得多.

\subsection{定义}\label{sec:galois-1-3-1}
设 $A$ 是一个环. 一个 $A$-\newterm{galois-module}{模}{module} $M$ 是配备了 $A$ 的一个``作用''的 Abel 群：
\begin{align*}
 A\times M&\longrightarrow M,\\
 (a,m)&\longmapsto a\cdot m.
\end{align*}
对所有 $a,a'\in A$ 和 $m,m'\in M$，这个作用满足下列公理：
\begin{enumerate}[label=\roman*)]
\item $a\cdot(m+m')=a\cdot m+a\cdot m'$（作用的线性性）.
\item $(a+a')\cdot m=a\cdot m+a'\cdot m$ 且 $(aa')\cdot m=a\cdot(a\cdot m)$.\footnote{译注：第二个等式右侧内层的标量原文写作 $a$，按通常的模公理应为 $a'$. 此处保留原式.}
\item $1\cdot m=m$.
\end{enumerate}
也称 $M$ 为``$A$ 上的模''.
\begin{remark*}
公理 i) 表明映射 $m\mapsto a\cdot m$ 是 Abel 群 $M$ 的一个自同态；公理 ii) 和 iii) 则表明，由此得到的映射 $A\longrightarrow\operatorname{End}(M)$ 是一个环同态（注意，$\operatorname{End}(M)$ 的乘法由自同态的复合给出，因此不交换，但这不影响同态的概念）. 反过来，每个环同态 $\psi\colon A\longrightarrow\operatorname{End}(M)$（这里取 Abel 群 $M$ 的自同态）都通过 $a\cdot m:=\psi(a)(m)$ 定义一个环结构.\footnote{译注：原文写作``环结构''（structure d'anneau），此处按上下文应为 $M$ 上的 $A$-模结构.}
\end{remark*}
我们常简化记号，将 $a\cdot m$ 写成 $am$. 下面的定义是自然的：
\begin{definition*}
$A$-模的\newterm{galois-module-morphism}{同态}{module homomorphism} $\varphi\colon M\longrightarrow M'$ 是一个 Abel 群同态，且它是 $A$-线性的，即对任意 $a\in A$ 和 $m\in M$，都有 $\varphi(am)=a\varphi(m)$.
\end{definition*}
我们用 $\operatorname{Hom}_A(M,M')$ 表示 $A$-模同态组成的集合，用 $\operatorname{End}_A(M)$ 表示 $A$-模 $M$ 的自同态组成的集合. 注意，它们本身也通过下列公式带有自然的 $A$-模结构：
\begin{align*}
 a\cdot\varphi&\colon m\longmapsto a\varphi(m),&
 \varphi+\varphi'&\colon m\longmapsto\varphi(m)+\varphi'(m).
\end{align*}
与往常一样，$A$-模同构是一个左右均可逆的同态. 容易验证，一个同态是同构，当且仅当它作为映射是双射.
\begin{example*}[{$A=\mathbb Z$}]
每个 Abel 群都具有唯一的 $\mathbb Z$-模结构，由唯一的环同态 $\mathbb Z\longrightarrow\operatorname{End}(M)$ 给出. 因此，$\mathbb Z$-模无非就是 Abel 群，而每个 Abel 群同态也都是 $\mathbb Z$-模同态.
\end{example*}
\begin{example*}[{$A$ 是域}]
域 $K$ 上的模就是 $K$-向量空间，而 $K$-模同态就是 $K$-线性映射.
\end{example*}

\subsection{标量限制}\label{sec:galois-1-3-2}
设 $\varphi\colon B\longrightarrow A$ 是环同态，$M$ 是 $A$-模. 复合映射
\[
\begin{tikzcd}
 B\arrow[r,"\varphi"] & A\arrow[r] & \operatorname{End}_{\mathbb Z}(M)
\end{tikzcd}
\]
赋予 $M$ 一个 $B$-模结构，由 $b\cdot m:=\varphi(b)m$ 给出. 称这个 $B$-模 $M$ 为 $A$-模 $M$``沿 $\varphi$ 的\newterm{galois-restriction-scalars}{标量限制}{restriction of scalars}''. 若要消除记号 $M$ 对所考虑的结构造成的歧义，可以将这个 $B$-模记作 $\varphi^*M$、$M_{|B}$ 或 $\operatorname{Res}_A^B(M)$.
\begin{example*}
$A$-模 $M$ 沿典范同态 $\mathbb Z\longrightarrow A$ 的标量限制就是 $M$ 的底层 Abel 群.
\end{example*}
注意，$A$-模同态 $M\longrightarrow N$ 也是 $B$-模同态 $\varphi^*M\longrightarrow\varphi^*N$，但反过来一般不成立. 因此有包含关系
\begin{align*}
 \operatorname{Hom}_A(M,N)\subset\operatorname{Hom}_B(M,N).
\end{align*}
\begin{exercise*}
验证，当 $\varphi$ 是满射时，上式为等式.
\end{exercise*}

\subsection{子模与商模}\label{sec:galois-1-3-3}
自然地，$M$ 的 $A$-\newterm{galois-submodule}{子模}{submodule}是 $M$ 的一个子群，并且在 $A$ 对 $M$ 的作用下稳定.
\begin{remark*}
乘法使 $A$ 成为一个 $A$-模. $A$ 的 $A$-子模无非就是 $A$ 的理想.
\end{remark*}
\begin{example*}
同态 $\varphi\colon M\longrightarrow M'$ 的像 $\operatorname{Im}(\varphi)$ 是 $M'$ 的 $A$-子模，而它的核 $\operatorname{Ker}(\varphi)=\varphi^{-1}(0)$ 是 $M$ 的 $A$-子模. 更一般地，若 $N$ 是 $M$ 的子模，则其像 $\varphi(N)$ 是 $M'$ 的子模；若 $N'$ 是 $M'$ 的子模，则其原像 $\varphi^{-1}(N')$ 是 $A$ 的子模.\footnote{译注：原文在此写作``$A$ 的子模''，按映射的定义域应为``$M$ 的子模''.}
\end{example*}
设 $M$ 是 $A$-模，$N$ 是 $M$ 的 $A$-子模. 考虑 $M$ 关于等价关系
\begin{align*}
 m\sim m'\quad\Longleftrightarrow\quad m-m'\in N
\end{align*}
的商集 $M/N$，并用 $\pi\colon M\longrightarrow M/N$ 表示典范投影. 与前面一样，我们视方便将 $m$ 的等价类记作 $\pi(m)$、$\overline m$、$m+N$ 或 $m\pmod N$.
\begin{proposition*}
$M/N$ 上存在唯一的 $A$-模结构，使得 $\pi$ 成为 $A$-模同态.
\end{proposition*}
\begin{proof}
论证与构造商环 $A/I$ 时相同. 唯一性来自 $\pi$ 的满射性，它迫使下列公式成立：$a\overline m=\overline{am}$ 和 $\overline m+\overline{m'}=\overline{m+m'}$. 为证明存在性，必须验证这些公式有意义，即
\begin{align*}
 m\sim m'&\ \Longrightarrow\ am\sim am',\\
 m\sim m_1,\quad m'\sim m'_1&\ \Longrightarrow\ (m+m')\sim(m_1+m'_1).
\end{align*}
这直接由 $N$ 是 $A$-子模得到.
\end{proof}
\begin{exercise*}
证明，$\pi^{-1}$ 诱导双射
\begin{align*}
 \{M/N\text{ 的子模}\}&\xrightarrow{\sim}\{M\text{ 中包含 }N\text{ 的子模}\},
\end{align*}
其逆双射为 $P\longmapsto\pi(P)=P/N$.
\end{exercise*}
与各种商对象一样，$M/N$ 也可以由泛性质刻画.
\begin{proposition*}
设 $\psi\colon M\longrightarrow M'$ 是满足 $\psi(N)=\{0\}$ 的 $A$-模同态. 存在唯一的 $A$-模同态 $\overline\psi\colon M/N\longrightarrow M'$，使得 $\psi=\overline\psi\circ\pi$.
\end{proposition*}
\begin{corollary*}
每个同态 $M\xrightarrow{\psi}M'$ 都具有唯一的分解
\[
\begin{tikzcd}
 M\arrow[r,two heads] & M/\operatorname{Ker}(\psi)\arrow[r,"\sim"] & \operatorname{Im}(\psi)\arrow[r,hook] & M'.
\end{tikzcd}
\]
\end{corollary*}
现在设 $P$ 是 $M$ 的一个包含 $N$ 的子模. 典范投影 $\pi_P\colon M\longrightarrow M/P$ 的核包含 $N$，因而上述命题给出 $\pi_P$ 的分解
\[
\begin{tikzcd}
 M\arrow[r,"\pi_N"] & M/N\arrow[r,"\overline\pi_P"] & M/P.
\end{tikzcd}
\]
\begin{corollary*}[{第一同构定理}]
$\overline\pi_P$ 诱导同构
\begin{align*}
 (M/N)/(P/N)\xrightarrow{\sim}M/N.
\end{align*}
\footnotemark
\end{corollary*}
\footnotetext{译注：原文在此同构及下文证明中的目标均写作 $M/N$，按 $\overline\pi_P$ 的定义应为 $M/P$. 此处保留原式.}
\begin{proof}
由于 $\pi_N$ 是满射，有 $\operatorname{Ker}(\overline\pi_P)=\pi_N(\operatorname{Ker}(\pi_P))=\pi_N(P)=P/N$，所以 $\overline\pi_P$ 通过单同态 $(M/N)/(P/N)\hookrightarrow M/N$ 分解. 这个同态也是满射，因为 $\pi_P$ 是满射，从而 $\overline\pi_P$ 也是满射.
\end{proof}
现在考虑如下情形：设 $M$ 是 $A$-模，$N,P$ 是 $M$ 的两个子模. 定义它们的和为
\begin{align*}
 N+P:=\{m\in M\mid\exists n\in N,\ \exists p\in P,\ m=n+p\}.
\end{align*}
这显然是 $M$ 的一个包含 $N$ 和 $P$ 的 $A$-子模（请验证！）. 同样，交集 $N\cap P$ 也是 $M$ 的 $A$-子模. 复合同态
\[
\begin{tikzcd}
 \rho\colon N\arrow[r,hook] & N+P\arrow[r,two heads] & (N+P)/P
\end{tikzcd}
\]
的核显然包含 $N\cap P$，所以这个同态通过同态
\begin{align*}
 N/(N\cap P)\longrightarrow(N+P)/P
\end{align*}
分解.
\begin{proposition*}[{第二同构定理}]
这个同态是同构：
\begin{align*}
 N/(N\cap P)\xrightarrow{\sim}(N+P)/P.
\end{align*}
\end{proposition*}
\begin{proof}
为证明单射性，需要证明 $\operatorname{Ker}(\rho)=N\cap P$，这是显然的. 为证明满射性，需要看到 $(N+P)/P$ 的每个元素都能沿投影 $N+P\twoheadrightarrow(N+P)/P$ 提升为 $N$ 中的一个元素；由商的定义，这也是直接的.
\end{proof}

\subsection{直和与直积}\label{sec:galois-1-3-4}
设 $I$ 是一个集合，$(M_i)_{i\in I}$ 是一族以 $I$ 为指标集的 $A$-模（可以考虑 $I=\mathbb N$ 或 $I=\{1,\ldots,r\}$）. 回顾，笛卡尔积 $\prod_{i\in I}M_i$ 是所有以 $I$ 为指标集、且对每个 $i$ 都有 $m_i\in M_i$ 的族 $(m_i)_{i\in I}$ 组成的集合. 用下式赋予这个笛卡尔积一个 $A$-模结构：
\begin{align*}
 (m_i)_{i\in I}+(m'_i)_{i\in I}&:=(m_i+m'_i)_{i\in I},\\
 a\cdot(m_i)_{i\in I}&:=(am_i)_{i\in I}.
\end{align*}
（留给读者验证：这确实是一个 $A$-模，其零元是族 $(0)_{i\in I}$.）这个模称为 $M_i$ 的\newterm{galois-product-module}{直积}{product of modules}，记作
\begin{align*}
 \prod_{i\in I}M_i,\qquad\text{当 }I=\{1,\ldots,r\}\text{ 时简记为 }M_1\times\cdots\times M_r.
\end{align*}
对每个 $j\in I$，投影 $(m_i)_{i\in I}\longmapsto m_j$ 是一个 $A$-模同态
\begin{align*}
 \prod_{i\in I}M_i\xrightarrow{\pi_j}M_j.
\end{align*}
现在定义
\begin{align*}
 \bigoplus_{i\in I}M_i:=\left\{(m_i)_{i\in I}\in\prod_{i\in I}M_i\ \middle|\ m_i=0\text{ 对几乎所有 }i\text{ 成立}\right\},
\end{align*}
即有限支撑族组成的子集（也就是说，除 $I$ 的一个有限子集外，都有 $m_i=0$）. 可见，它是 $\prod_{i\in I}M_i$ 的一个 $A$-子模，称为 $M_i$ 的\newterm{galois-direct-sum}{直和}{direct sum}或\newterm{galois-coproduct}{余积}{coproduct}.
\begin{remark*}
若 $I$ 有限，则 $\bigoplus_{i\in I}M_i=\prod_{i\in I}M_i$. 当 $I=\{1,\ldots,r\}$ 时，也记作 $\bigoplus_{i=1}^rM_i$ 或简记作 $M_1\oplus\cdots\oplus M_r$.
\end{remark*}
对每个 $j\in I$，有映射
\begin{align*}
 \iota_j\colon M_j\hookrightarrow\bigoplus_{i\in I}M_i,
\end{align*}
将 $m$ 映为满足 $i\ne j$ 时 $m_i=0$、$m_j=m$ 的族 $(m_i)_{i\in I}$. 这个映射显然是 $A$-模同态.

这些模连同各自的同态族分别满足一个泛性质，且这两个性质在某种意义上互为``对偶''.
\begin{proposition*}
\emergencystretch=3em
i) 设 $N$ 是 $A$-模，并且对每个 $i\in I$ 都给定同态 $\psi_i\colon N\longrightarrow M_i$. 则存在唯一的同态 $\Psi\colon N\longrightarrow\prod_{i\in I}M_i$，使得对所有 $i$ 都有 $\psi_i=\pi_i\circ\Psi$. 换言之，映射 $\Psi\longmapsto(\pi_i\circ\Psi)_{i\in I}$ 诱导双射
\begin{align*}
 \operatorname{Hom}_A\left(N,\prod_{i\in I}M_i\right)\xrightarrow{\sim}\prod_{i\in I}\operatorname{Hom}_A(N,M_i).
\end{align*}
ii) 设 $N$ 是 $A$-模，并且对每个 $i\in I$ 都给定同态 $\psi_i\colon M_i\longrightarrow N$. 则存在唯一的同态 $\Psi\colon\bigoplus_{i\in I}M_i\longrightarrow N$，使得对所有 $i$ 都有 $\psi_i=\Psi\circ\iota_i$. 换言之，映射 $\Psi\longmapsto(\Psi\circ\iota_i)_{i\in I}$ 诱导双射
\begin{align*}
 \operatorname{Hom}_A\left(\bigoplus_{i\in I}M_i,N\right)\xrightarrow{\sim}\prod_{i\in I}\operatorname{Hom}_A(M_i,N).
\end{align*}
\end{proposition*}
\begin{proof}
这些全都是形式上的论证. i) 对于存在性，只需令 $\Psi(n):=(\psi_i(n))_{i\in I}$. 对于唯一性，若 $\Psi'$ 是另一个同态，则对所有 $n$，$\Psi(n)-\Psi'(n)$ 在每个投影 $\pi_i$ 下的像都为零，因此它本身为零.

 ii) 对于存在性，只需令 $\Psi((m_i)_{i\in I}):=\sum_{i\in I}\psi_i(m_i)$；这是有意义的，因为族 $\psi_i(m_i)$ 几乎处处为零（这个和中只有有限多个非零项）. 对于唯一性，若 $\Psi'$ 是另一个解，则 $\Psi-\Psi'$ 在每个 $\iota_i$ 的像 $\iota_i(M_i)$ 上都为零. 而 $\bigoplus_{i\in I}M_i$ 的每个元素都是这种形式的元素之和（若 $I$ 无限，这对 $\prod_{i\in I}M_i$ 的元素并不成立）.
\end{proof}




\subsection{子模的和与生成的模}\label{sec:galois-1-3-5}
设 $M$ 是一个 $A$-模. 两个子模的交仍是子模，因此可以谈论包含给定子集 $E\subset M$ 的``最小子模'' $M(E)$. 它也是所有包含 $E$ 的子模的交，称为 $E$ 生成的子模. 具体地，它是集合
\begin{align*}
M(E)=\{m=a_1e_1+\cdots+a_re_r,\ r\in\mathbb N,\ e_i\in E,\ a_i\in A\}.
\end{align*}
现在设给定一族以集合 $I$ 为指标的子模 $(M_i)_{i\in I}$. 用
\begin{align*}
\sum_{i\in I}M_i,\qquad\text{或在 $I=\{1,\ldots,r\}$ 时简记为 }M_1+\cdots+M_r,
\end{align*}
表示并集 $\bigcup_{i\in I}M_i$ 生成的 $M$ 的子模. 具体地，
\begin{align*}
\sum_{i\in I}M_i=\left\{\sum_{i\in J}m_i,\ J\subset I\text{ 有限},\ m_i\in M_i\right\}.
\end{align*}
更形式地，考虑由余积的泛性质给出的、与包含映射 $\psi_i\colon M_i\hookrightarrow M$ 对应的态射 $\Psi\colon\bigoplus_{i\in I}M_i\longrightarrow M$. 则
\begin{align*}
\sum_{i\in I}M_i=\operatorname{Im}(\Psi).
\end{align*}
\begin{definition*}
若上述态射 $\Psi$ 是单射，就称子模 $M_i$ 的和是直和. 此时它诱导同构 $\bigoplus_{i\in I}M_i\xrightarrow{\sim}\sum_{i\in I}M_i$.
\end{definition*}
$\Psi$ 的单射性也可以表述为：对任意有限族元素 $m_i\in M_i$，有 $\sum_{i\in I}m_i=0\Rightarrow\forall i\in I,\ m_i=0$.
\begin{exercise*}
证明：子模 $M_1,\ldots,M_r$ 的和是直和，当且仅当对每个 $i=1,\ldots,r$，都有 $M_i\cap\sum_{j\ne i}M_j=\{0\}$.
\end{exercise*}
\begin{remark*}
当 $A=k$ 是域时，我们知道向量空间 $V$ 的任意子空间 $W$ 都有补空间，即存在 $V$ 的子空间 $W'$ 使 $V=W\oplus W'$. 一般而言，这不再成立. 例如，当 $A=\mathbb Z$ 时，$M=\mathbb Z$ 的子模 $2\mathbb Z$ 没有补模，因为 $\mathbb Z$ 的每个非零子模与 $2\mathbb Z$ 的交均非零.
\end{remark*}
若模 $M$ 的子模 $N$ 有补模，就称它为 $M$ 的\newterm{galois-direct-summand}{直和因子}{direct summand}.

\subsection{自由模}\label{sec:galois-1-3-6}
直和的一个重要特例是对每个 $i\in I$ 都有 $M_i=A$ 的情形. 此时记
\begin{align*}
A^I:=\prod_{i\in I}A\qquad\text{和}\qquad A^{(I)}:=\bigoplus_{i\in I}A.
\end{align*}
当 $I$ 有限时，当然有 $A^{(I)}=A^I$，通常采用更简单的后一种记号. 当 $I=\{1,\ldots,r\}$ 时，也简记为 $A^r$ 或 $A^{\oplus r}$，而不写 $A^{\{1,\ldots,r\}}$. 对 $i\in I$，用 $e_i$ 表示 $A^{(I)}$ 中除第 $i$ 个分量为 $1$ 外其余分量均为零的元素. 例如，若 $I=\{1,\ldots,r\}$，则 $e_i=(0,\ldots,1,\ldots,0)$，其中 $1$ 位于第 $i$ 个位置.
\begin{proposition*}
\begin{enumerate}[label=\roman*)]
\item $A^{(I)}$ 的每个元素都唯一地写成 $\sum_{i\in I}a_ie_i$，其中 $(a_i)_{i\in I}$ 是 $A$ 中一个除有限项外均为零的元素族.
\item $A$-模 $A^{(I)}$ 具有如下泛性质：对任意 $A$-模 $M$ 和 $M$ 中任意元素族 $(m_i)_{i\in I}$，都存在唯一的 $A$-模态射 $\Psi\colon A^{(I)}\longrightarrow M$ 将 $e_i$ 映为 $m_i$. 换言之，映射 $\Psi\mapsto(\Psi(e_i))_{i\in I}$ 是一个双射
\begin{align*}
\operatorname{Hom}_A(A^{(I)},M)\xrightarrow{\sim}M^I.
\end{align*}
\end{enumerate}
\end{proposition*}
\begin{proof}
i) 由构造可得.

ii) 首先注意，对模 $M$ 的任意元素 $m$，存在唯一的 $A$-模态射 $A\xrightarrow{\psi_m}M$ 将 $1$ 映为 $m$. 它由 $\varphi_m(a):=am$ 定义. 因此，对每个 $i$ 都有一个态射 $\psi_{m_i}\colon A\longrightarrow M$，再使用直和的泛性质即可. 具体地，就是 $\Psi(\sum_{i\in I}a_ie_i)=\sum_{i\in I}a_im_i$.
\end{proof}
\begin{definition*}
设 $(m_i)_{i\in I}$ 是 $M$ 中的一族元素，$\Psi\colon A^{(I)}\longrightarrow M$ 是相应的态射. 称这个元素族为：
\begin{itemize}
\item \newterm{galois-linearly-independent-family}{线性无关族}{linearly independent family}，若 $\Psi$ 是单射；这等价于条件 $\sum_i a_im_i=0\Rightarrow\forall i,\ a_i=0$.
\item \newterm{galois-generating-family}{生成族}{generating family}，若 $\Psi$ 是满射；这等价于每个 $m\in M$ 都可以写成 $\sum_{i\in I}a_im_i$，其中对每个 $i$ 都有 $a_i\in A$.
\item $M$ 的一个\newterm{galois-basis}{基}{basis}，若 $\Psi$ 是同构；这等价于每个 $m\in M$ 都唯一地写成 $\sum_{i\in I}a_im_i$，其中对每个 $i$ 都有 $a_i\in A$.
\end{itemize}
若 $M$ 有一个有限生成族，就称它是\newterm{galois-finitely-generated-module}{有限生成模}{finitely generated module}.
\end{definition*}
\begin{example*}
设 $A=\mathbb Z$，$M=\mathbb Z$.
\begin{itemize}
\item 元素族 $\{2,3\}$ 生成 $M$，因为由 Bézout 等式，每个 $n\in\mathbb Z$ 都可写为 $2a+3b$. 但它不是基，因为 $0=2\cdot3-3\cdot2$.
\item 元素族 $\{2\}$ 线性无关，但不是生成族，因此不是基.
\item $M$ 的基只有 $\{1\}$ 和 $\{-1\}$.
\end{itemize}
\end{example*}
\begin{example*}
更一般地，对 $M=A$，任何包含两个不同元素 $a,a'$ 的元素族都不是线性无关族，因为有关系 $a.a'-a'.a=0$. 因此，线性无关族是单元素集 $\{a\}$，其中 $a$ 是 $A$ 的正则元. 此外，这样的单元素集是基，当且仅当 $a$ 可逆.
\end{example*}
\begin{definition*}
若 $A$-模 $M$ 有一个基，就称它是\newterm{galois-free-module}{自由模}{free module}. 每个基的选择都会诱导同构 $A^{(I)}\xrightarrow{\sim}M$，其中 $I$ 是适当的集合.
\end{definition*}
\begin{example*}
下面是几个非自由模的例子.
\begin{itemize}
\item 若 $I$ 是 $A$ 的非零真理想，则 $A$-模 $M:=A/I$ 不是自由模. 事实上，$M$ 没有任何线性无关族，因为每个 $i\in I\setminus\{0\}$ 的作用都将 $M$ 消为零.
\item 设 $A=\mathbb C[X,Y]$，$M=(X,Y)$（由 $X$ 和 $Y$ 生成的理想）. 因为 $M\subset A$，由前例可知，$M$ 的线性无关族只有形如 $\{f(X,Y)\}$ 的单元素集，其中 $f(X,Y)$ 是常数项为零的非零多项式. 但容易看出，单元素集绝不可能生成 $M$.
\end{itemize}
\end{example*}
自由模具有域上线性空间的一些良好性质. 例如，若 $M$ 是以 $(e_i)_{i=1,\ldots,n}$ 为基的自由模，$N$ 是以 $(f_j)_{i=1,\ldots,m}$ 为基的自由模，$\varphi$ 是 $A$-模态射 $N\longrightarrow M$，就可以写 $\varphi(f_j)=\sum_{i=1}^n a_{ij}e_i$. 由此得到与 $A$ 上的 $n\times m$ 矩阵之间的双射（甚至是 $A$-模同构）
\begin{align*}
\operatorname{Hom}_A(N,M)\xrightarrow{\sim}\mathcal M_{n,m}(A).
\end{align*}
当 $N=M$ 时，这甚至是一个环同构.

不过，仍有一些显著区别. 下面举几个例子.
\begin{itemize}
\item 有限秩自由模的单射自同态未必是满射. 例如 $A=\mathbb Z=M$，$\varphi$ 是 $M$ 的自同态 $m\mapsto2m$.
\item 未必能从生成族中抽出一个基. 例如 $A=\mathbb Z=M$，生成族为 $\{2,3\}$.
\item 线性无关族未必能扩充为基. 仍用同一例子，取线性无关族 $\{2\}$.
\item 秩为 $n$ 的自由模的子模未必自由，也未必能由基数小于 $n$ 的元素族生成. 例如 $A=\mathbb C[X,Y]$，$n=1$，$M=AX+AY$.
\end{itemize}
不过，还是有一条好消息.
\begin{proposition*}
设 $M$ 是有限生成自由 $A$-模. 则它的所有基都是有限的，且基数相同. 这个基数称为 $M$ 的\newterm{galois-rank}{秩}{rank}. 此外，线性无关族（相应地，生成族）的基数不大于（相应地，不小于）$M$ 的秩.
\end{proposition*}
\begin{proof}
首先假设 $M$ 有一个有限基 $E=(e_1,\ldots,e_n)$，并设 $F=(f_1,\ldots,f_m)$ 是 $M$ 中的一族元素. 我们要证明：若 $F$ 是生成族，则 $m\geqslant n$. 由对称性可知，若 $F$ 是基，则 $n=m$.

为证明 $m\geqslant n$，将问题归结到基环是域的已知情形. 写 $f_j=\sum_{i,j}a_{ij}e_i$. 矩阵 $P=(a_{ij})_{i,j}\in\mathcal M_{n\times M}(A)$ 是一个态射 $A^m\xrightarrow{\psi}A^n$ 的矩阵，且 $F$ 是生成族，当且仅当 $\psi$ 是满射. 现在选择 $A$ 的一个极大理想 $\mathfrak m$（Zorn 引理保证存在按包含关系极大的真理想，因为真理想的任意递增并仍是真理想），并用 $\pi\colon A\longrightarrow A/\mathfrak m$ 表示商映射. 矩阵 $(\pi(a_{ij}))_{i,j}\in\mathcal M_{n\times m}(A/\mathfrak m)$ 是态射 $\overline\psi\colon(A/\mathfrak m)^m\longrightarrow(A/\mathfrak m)^n$ 的矩阵，并构成交换图
\[
\begin{tikzcd}
A^m \arrow[r,"\psi"] \arrow[d,"\pi^m"'] & A^n \arrow[d,"\pi^n"] \\
(A/\mathfrak m)^m \arrow[r,"\overline\psi"] & (A/\mathfrak m)^n
\end{tikzcd}
\]
（回顾，图的交换性是指 $\pi^n\circ\psi=\overline\psi\circ\pi^m$）. 因为 $\pi^n$ 是满射，可见
\begin{align*}
\psi\text{ 满射}&\Rightarrow\pi^n\circ\psi\text{ 满射}\\
&\Leftrightarrow\overline\psi\circ\pi^m\text{ 满射}\Rightarrow\overline\psi\text{ 满射}.
\end{align*}
又因 $A/\mathfrak m$ 是域，已知 $\overline\psi$ 满射 $\Rightarrow m\geqslant n$.

现在证明，若 $m>n$，则元素族 $F$ 线性相关. 设 $r$ 是使 $P$ 存在非零 $r\times r$ 子式的最大整数. 若 $r=0$，则所有 $f_i$ 均为零，$F$ 显然线性相关，因此假设 $r\geqslant1$. 还有 $r\leqslant n<m$. 适当重编号后，可设子式 $\mu_1:=\det((a_{ij})_{1\leqslant i\leqslant r,\,2\leqslant j\leqslant r+1})$ 非零. 对 $k=2,\ldots,r+1$，用 $\mu_k$ 表示子式 $\det((a_{ij})_{1\leqslant i\leqslant r,\,1\leqslant j\leqslant r+1,\,j\ne k})$. 此时，对每个 $i=1,\ldots,n$，和 $\sum_{k=1}^{r+1}(-1)^{k+1}a_{ik}\mu_k$ 或者是 $P$ 的一个 $r+1$ 阶子式，或者是有两行相同的矩阵的行列式，因此为零. 因为 $E$ 是基，所以 $\sum_{k=1}^{r+1}\mu_kf_k=0$；又因 $\mu_1\ne0$，可见 $F$ 线性相关.\footnote{译注：原文的最后一个线性关系未写交替符号；按前面的余子式展开，应为 $\sum_{k=1}^{r+1}(-1)^{k+1}\mu_kf_k=0$. 正文保留原式.}

最后，还要去掉最初关于 $M$ 有有限基的假设. 命题陈述中没有这个假设，只假设对某个集合 $I$ 有 $M\simeq A^{(I)}$. 因此假设 $I$ 无限，而 $M$ 另有一个有限生成族，设其基数为 $m$. 我们要导出矛盾. 为此，取一个基数为 $n>m$ 的集合 $J\subset I$. 有典范态射 $A^{(I)}\xrightarrow{\rho}A^J$ 投影到以 $J$ 为指标的分量. 这个态射是满射，因此元素族 $\rho(F)$ 生成 $A^J=A^n$；这与前面的讨论矛盾，因为 $m<n$.
\end{proof}


\subsection{Noether 模与 Noether 环}\label{sec:galois-1-3-7}
有限生成模具有下列性质（作为练习）：
\begin{itemize}
\item $M$ 有限生成且 $N\subset M$，则 $M/N$ 有限生成.
\item $N$ 有限生成且 $M/N$ 有限生成，则 $M$ 有限生成.
\end{itemize}
然而，有限生成模的子模未必有限生成！下面给出两个例子.
\begin{example*}
令 $M=A=\mathbb Z[X_i,\ i\in\mathbb N]$（无限多个变量的多项式环，见下文），$N$ 为所有常数项为零的多项式组成的理想. $N$ 中的任意有限元素族都包含在 $X_1,\ldots,X_n$ 生成的理想中，只要 $n$ 足够大；而这个理想不包含 $X_{n+1}$. 因此 $N$ 不有限生成.
\end{example*}
\begin{example*}
令 $M=A=\overline{\mathbb Z}$，$N$ 为序列 $(x_n)_{n\in\mathbb N}$ 生成的理想，其中 $x_0=2$，$x_n$ 是 $x_{n-1}$ 的一个平方根. 同样，$N$ 的任意有限生成子模都包含在 $x_1,\ldots,x_n$ 生成的子环中，只要 $n$ 足够大，而后者不包含 $x_{n+1}$.\footnote{译注：原文这里写的是``子环''，并据此断言任意有限生成子模都包含于其中. 对 $\overline{\mathbb Z}$-子模，这一包含关系一般不成立；正文保留原文论证.}
\end{example*}
为避免这些病态现象，引入 Noether 模与 Noether 环的概念（以提出这些概念的 Emmy Noether 命名）.
\begin{proposition*}
设 $M$ 是一个 $A$-模. 下列性质等价：
\begin{enumerate}[label=\roman*)]
\item $M$ 的每个 $A$-子模都有限生成.
\item 每个递增的 $A$-子模序列都从某一项起稳定.
\item $M$ 的 $A$-子模组成的任意非空集合，都有按包含关系的极大元.
\end{enumerate}
\end{proposition*}
\begin{proof}
i) $\Rightarrow$ ii). 设 $(M_n)_{n\in\mathbb N}$ 是递增子模序列. 则并集 $M=\bigcup_{n\in\mathbb N}M_n$ 也是子模（请验证！）.\footnote{译注：原文在这里再次使用 $M$ 表示该并集；它不一定等于命题中给定的整个模.} 由性质 i)，它由有限个元素生成；当 $n$ 足够大时，这些元素全都属于某个 $M_n$. 因此 $M=M_n$，且对所有 $N\geqslant n$ 都有 $M_N=M_n$.

ii) $\Rightarrow$ iii). 证明其逆否命题. 假设存在一个由 $M$ 的 $A$-子模组成的集合，其中没有极大元. 于是可以递归构造严格递增序列，它永远不会稳定.

iii) $\Rightarrow$ i). 设 $M'$ 是 $M$ 的子模. 考虑 $M'$ 的所有有限生成子模组成的集合，它非空，因为包含 $\{0\}$. 由性质 iii)，它有一个极大元 $N$. 任取 $m'\in M'$. $M'$ 的子模 $N+(m')$ 有限生成，因而由 $N$ 的极大性，它包含于 $N$. 故 $m'\in N$，于是 $M'=N$ 有限生成.
\end{proof}
\begin{definition*}
满足上个命题中等价条件的 $A$-模 $M$ 称为\newterm{galois-noetherian-module}{Noether 模}{Noetherian module}. 若环 $A$ 作为 $A$-模是 Noether 模，就称 $A$ 为\newterm{galois-noetherian-ring}{Noether 环}{Noetherian ring}.
\end{definition*}
因此，若一个环的所有理想都有限生成，则它是 Noether 环. 特别地，主理想环（即所有理想都是主理想的环）是 Noether 环.
\begin{proposition*}
\begin{enumerate}[label=\roman*)]
\item 设 $M$ 是 $A$-模，$N$ 是其子模. 则 $M$ 是 Noether 模，当且仅当 $N$ 和 $M/N$ 都是 Noether 模.
\item 有限个模的直和是 Noether 模，当且仅当每个直和因子都是 Noether 模.
\item Noether 环上的任意有限生成模都是 Noether 模.
\item 设 $B\xrightarrow{\varphi}A$ 是环态射. 若 $M$ 作为 $B$-模是有限生成模（相应地，Noether 模），则它作为 $A$-模也是有限生成模（相应地，Noether 模）.
\end{enumerate}
\end{proposition*}
\begin{proof}
i) 假设 $M$ 是 Noether 模. 则 $N$ 的每个子模都是 $M$ 的子模，因而有限生成，所以 $N$ 是 Noether 模. 此外，$M/N$ 的每个子模 $\overline P$ 都是 $M$ 的某个子模 $P$（包含 $N$）在投影 $M\longrightarrow M/N$ 下的像，而 $P$ 由有限个元素 $m_1,\ldots,m_r$ 生成. 因此 $\overline P$ 由 $\overline m_1,\ldots,\overline m_r$ 生成，故 $M/N$ 是 Noether 模.

现在假设 $N$ 和 $M/N$ 都是 Noether 模，并设 $P$ 是 $M$ 的子模. 则 $P\cap N$ 是 $N$ 的子模，因而有限生成. 另一方面，商模 $P/(P\cap N)$ 典范同构于商模 $(N+P)/N$（第二同构定理），后者是 $M/N$ 的子模（见第~\ref{sec:galois-1-3-3}~分目的练习），因而有限生成. 由此可知 $P$ 本身有限生成：若 $P\cap N$ 由 $p_1,\ldots,p_r$ 生成，$P/(N\cap P)$ 由 $\overline p_{r+1},\ldots,\overline p_s$ 生成，则 $P$ 由 $p_1,\ldots,p_s$ 生成，其中 $p_{r+1},\ldots,p_s$ 分别是 $\overline p_{r+1},\ldots,\overline p_s$ 在 $P$ 中任取的提升.

ii) 有两个直和因子时，这是 i) 的特例，因为若 $M=M_1\oplus M_2$，则到 $M_2$ 的投影诱导同构 $M/M_1\xrightarrow{\sim}M_2$. 对 $n$ 个因子的情形直接归纳即可.

iii) 假设 $M$ 由 $m_1,\ldots,m_r$ 生成. 与这个元素族对应的是一个满态射 $A^r\longrightarrow M$（将 $e_i$ 映为 $m_i$）. 因此 $M$ 是 $A^r$ 的商模；由 ii)，$A^r$ 是 Noether 模，再由 i)，$M$ 是 Noether 模.

iv) 对``有限生成''，只需注意，任意 $B$-生成族当然也是 $A$-生成族. 对``Noether''，只需注意，$M$ 的递增 $A$-子模序列也是递增 $B$-子模序列.
\end{proof}
\begin{corollary*}
Noether 环的任意商环都是 Noether 环.
\end{corollary*}
\begin{proof}
设 $A=B/J$，其中 $B$ 是 Noether 环. $B$-模 $B/J$ 有限生成（甚至是单生成的），因而是 Noether 模；由前一命题的 iv)，它作为 $A$-模也是 Noether 模. 因此环 $A$ 是 Noether 环.
\end{proof}
Noether 环的一个优点是，它们具有足够多的不可约元（例如，与之相反，环 $\overline{\mathbb Z}$ 没有任何不可约元）.
\begin{theorem*}
设 $A$ 是 Noether 整环. 每个非零且不可逆的元素都是不可约元的乘积.
\end{theorem*}
\begin{proof}
考虑集合 $\mathcal I$，它由所有形如 $(a)$ 的主理想组成，其中 $a$ 非零、不可逆，且不是不可约元的乘积. 若 $\mathcal I$ 非空，由于假设 $A$ 是 Noether 环，它有一个极大元 $(a)$. 因为 $a$ 不可约性不成立，可以写 $a=bc$，其中 $b,c$ 均不可逆. 此时 $(b)$ 和 $(c)$ 都严格包含 $(a)$. 事实上，若例如 $(b)=(a)$，即对某个 $d\in A$ 有 $b=ad$，则 $a=acd$，从而 $cd=1$（$A$ 是整环），这与 $c$ 不可逆矛盾. 现在，由于 $a$ 不是不可约元的乘积，$b$ 或 $c$ 也不是不可约元的乘积，但这与选择 $(a)$ 为 $\mathcal I$ 的极大元矛盾.
\end{proof}
\begin{remark*}
不妨明确说明整环中不可约元与主理想之间的关系. 首先注意，元素 $a\in A$ 不可约，当且仅当主理想 $(a)$ 在 $A$ 的所有真主理想中极大. 事实上，设 $a=bc$，其中 $b$ 不可逆. 则 $(b)$ 是包含 $(a)$ 的真主理想，且可见 $(b)=(a)$ 当且仅当 $c$ 可逆（注意，$b=ad\Rightarrow a=acd\Rightarrow1=cd$，因为假设 $A$ 是整环）. 事实上，在任意整环中，映射 $a\mapsto(a)$ 诱导双射
\begin{align*}
&\{\text{按相伴关系取等价类的不可约元}\}\\
&\qquad\longleftrightarrow\{\text{在真主理想中极大的主理想}\}.
\end{align*}
\end{remark*}

